\documentclass[../../../include/open-logic-section]{subfiles}

\begin{document}

\olfileid{sth}{story}{rus}
\olsection{د رسل تناقض (بيا)}

په \olref[sfr][][]{part} کښې مو له ساده سټ تيورۍ سره کار
کړے و۔ خو د يو \emph{ډېر} ساده تصور له مخې، سټونه يوازې
د پريديکاتونو د صدق سټونه دي۔ دا ساده فکر به دا اصل لازمي کړي:

\begin{defish}
  \emph{ساده ادراک۔} د هر فارمول $\phi$ لپاره $\Setabs{x}{\phi(x)}$ وجود لري۔
\end{defish}

که څه هم دا اصل زړه راکښونکے دے، په ثبوت سره ناسازګاره دے۔
دا مو په \olref[sfr][set][rus]{sec} کښې ليدلي وو، خو نتيجه
دومره مهمه او دومره نېغه ده چې بيا ويل يې ګټه لري۔ ټکے په ټکے۔

\begin{thm}[د رسل تناقض]
سټ $R = \Setabs{x}{x \notin x}$ نشت۔
\end{thm}

\begin{proof}
که $R = \Setabs{x}{x \notin x}$ وجود ولري، نو
$R \in R$ که او يوازې که $R \notin R$، چې دا تناقض دے۔
\end{proof}

رسل دا نتيجه د 1901 په جون کښې وموندله۔ (خو تناقض يې په
کټ مټ هغې بڼې کښې نۀ و وړاندې کړے چې اوس مو وړاندې کړ؛ ځکه
هغه د فرېګه سټ تيوري څېړله، چې په \emph{Grundgesetze} کښې
بيان شوې ده۔ دې ته به په \olref[blv]{sec} کښې بيا راوګرځو۔)
رسل د 1902 د جون په 16مه فرېګه ته وليکل او د فرېګه د نظام
ناسازګاري يې ورته تشريح کړه۔ د خطونو او يو څه شاليد لپاره
\citet[pp.~124--8]{Heijenoort1967} وګورئ۔

ټينګار پکار دے چې دا د دوو کرښو ثبوت د \emph{خالص منطق}
يوه نتيجه ده۔ دا منو چې په نښه $\Setabs{x}{x \notin x}$ کښې
مو په پټه يو (غير منطقي؟)\ اصل، د غړو له مخې برابري، کارولے
دے؛ ځکه $\Setabs{x}{\phi(x)}$ د $\phi$ شرط پوره کوونکو شيانو د \emph{هماغه يوازيني}
سټ لپاره دے (د غړو له مخې د برابرۍ په وجه)، که وجود ولري۔ خو
د غړو له مخې د برابرۍ ان لږه اشاره هم لرې کولے شو، که نتيجه
داسې بيان کړو: \emph{داسې سټ نشت چې غړي يې کټ مټ هغه سټونه
وي چې خپل ځان د غړي په توګه نۀ لري}۔ او دا له سټونو سره ډېره
ځانګړې اړيکه نۀ لري۔ لکه رسل چې پخپله وليدل، کټ مټ همداسې
استدلال به دې نتيجې ته مو ورسوي: \emph{هيڅ سړے کټ مټ د هغو
سړو ږيرې نۀ خريي چې خپلې ږيرې پخپله نۀ خريي}۔ يا:
\emph{د پګ نسل هيڅ سپے کټ مټ هغه پګ سپي نۀ بوي کوي چې خپل
ځانونه نۀ بوي کوي}۔ او داسې نور۔ د شېما په بڼه د نتيجې جوړښت
يوازې دا دے:
\[
\lnot \exists x \forall z(Rzx \liff \lnot R zz).
\]
او دا يوازې د لومړۍ درجې منطق يوه قضيه (د قضيو شېما) ده۔ نو
د خپلې سټ تيورۍ په لږ بدلولو د رسل له تناقضه نۀ شو بچ کېدے؛
دا تر هغې مخکې راولاړېږي چې سټ تيورۍ ته \emph{ورسېږو}۔ که
(کلاسيکي) د لومړۍ درجې منطق کاروو، نو دا خبره بايد
\emph{ومنو} چې سټ $R = \Setabs{x}{x\notin x}$ نشت۔

پايله دا ده: که ساده ادراک منل غواړئ او له ناسازګارۍ څخه
\emph{بچ کېدل} هم غواړئ، نو يوازې په \emph{سټ تيوري} کښې
لږ بدلون بس نۀ دے۔ بلکې خپل \emph{منطق} به له بنسټه بيا جوړوئ۔

البته، له غير کلاسيکي منطقونو سره سټ تيورۍ وړاندې شوې دي۔
خو، په لږو ټکو، معياري نۀ دي۔ د رسل د تناقض معياري چلند دا
دے چې د نۀ شتون د نېغ ثبوت په توګه يې ومنو او بيا زده کړو
چې څنګه ورسره کار پر مخ يوسو۔ مونږ به همدا لاره ونيسو۔

\end{document}
